\chapter{Тензорный квадрат Real-пары и термическое переоткрытие перехода}

Предыдущая проверка оставила два возможных источника недостающего множителя
два: тензорный квадрат обменной Real-пары и двойной $SO(3)$-казимир.
Настоящий аудит сначала закрывает вопрос статической кратности, а затем
проверяет более физичную альтернативу: не обязан ли канонический внешний
квадрат запускать переход при охлаждении, без усиления самого оператора.

\section{Внешний квадрат полной Real-пары}

Пусть одночастичное пространство имеет разложение
\begin{equation}
 H=V_+\oplus V_-,
 \qquad \dim V_+=\dim V_-=3,
\end{equation}
а обменный оператор равен
\begin{equation}
 \Phi_B=
 \begin{pmatrix}0&B^T\\B&0\end{pmatrix}.
\end{equation}
Двухчастичное фермионное пространство канонически раскладывается как
\begin{equation}
 \Lambda^2H
 =\Lambda^2V_+
 \oplus(V_+\otimes V_-)
 \oplus\Lambda^2V_-.
 \label{eq:v6-real-pair-exterior-decomposition}
\end{equation}
Это стандартная внешняя степень прямой суммы и обычная конструкция
фермионного Фок-пространства; см.
\path{arXiv:quant-ph/0511002}, \path{arXiv:1308.3275}.

Пусть $X=B^TB$ и $T=\Tr X$. Прямой compound-matrix расчёт даёт
\begin{align}
 \|\Lambda^2\Phi_B\|_F^2
 &=2e_2(X)+T^2,\
 \|\Lambda^2\Phi_B\|^2_{\rm same}
 &=2e_2(X),\
 \|\Lambda^2\Phi_B\|^2_{\rm mixed}
 &=T^2.
 \label{eq:v6-real-pair-exterior-norm-split}
\end{align}
Первое слагаемое приходит от двух сопряжённых одноориентационный переходов,
а смешанный сектор после нормировки $R=X/T$ превращается в константу единица.
Следовательно, он не усиливает рангово-чувствительный член.

\begin{theorem}[Тензорный квадрат не удваивает селектор]
Полная Real-пара уже содержит прямой и сопряжённый внешние пути. После
физического полуслед их суммарный рангово-чувствительный вес остаётся
$2e_2(R)$. Смешанный сектор даёт только нормировочную константу и не
создаёт недостающий множитель два.
\end{theorem}

\section{Ненормированный антисимметризатор не является проектором}

На $V\otimes V$ обозначим перестановку факторов через $\Sigma$. Тогда
ортогональный проектор на антисимметричную часть равен
\begin{equation}
 P_- =\frac{I-\Sigma}{2}.
\end{equation}
Для $B\otimes B$ выполнены точные тождества
\begin{align}
 \|(B\otimes B)P_-\|_F^2&=e_2(B^TB),\
 \|(B\otimes B)(I-\Sigma)\|_F^2&=4e_2(B^TB).
 \label{eq:v6-raw-versus-projected-exchange}
\end{align}
Значит, ненормированный коэффициент четыре возникает потому, что $I-\Sigma=2P_-$.
Он считает нормированный антисимметричный канал с амплитудой два.

Условное ожидание на обменно-инвариантную алгебру всегда задаёт
ортогональный проектор, а не ненормированную разность. Как и в предыдущих
проверках относительного родительского действия в Томе IV, нормировка
спектрального разрыва обязана удалять этот квадрат двойки. Поэтому
ненормированную норму нельзя считать
автоматическим следствием неразличимости частиц.

\begin{nogocriterion}[Запрет обменного двойного счёта]
Два порядка обмена одинаковых фермионных факторов не являются двумя
независимыми физическими путями. После антисимметричной факторизации они дают
один нормированный сектор. Использование $I-\Sigma$ вместо $P_-$ без
отдельного оператора является скрытым множителем четыре.
\end{nogocriterion}

\section{Точный, но пока недоступный казимир четыре}

Для векторного представления $SO(3)$ при стандартной нормировке
\begin{equation}
 -\sum_{a=1}^3J_a^2=2I_3.
\end{equation}
Внешний канал несёт векторный индекс слева и справа, поэтому
\begin{equation}
 C_2^{L}+C_2^{R}=2+2=4.
 \label{eq:v6-left-right-so3-casimir-four}
\end{equation}
Число в точности совпадает с ненормированным коэффициентом разности
прямого и перекрёстного каналов.

Однако проверка
\path{version5_spatial_so3_superconnection_parent_trace_gate.tex}
уже доказал, что конечный след $M_{35}$ не выбирает
$\mathfrak{so}(3)$-подалгебру и не создаёт её лапласиан. В текущем
обменном родителе также нет независимых левой и правой пространственных
$SO(3)$-связностей. Следовательно, совпадение казимира является сильной
структурной подсказкой, но не выведенным весом действия.

\section{Охлаждение вместо усиления оператора}

Статический аудит до сих пор фиксировал отношение энергии к энтропии
равным единице. Для настоящего термодинамического функционала энергия
умножается на обратную температуру. Канонический внешний квадрат даёт
\begin{equation}
 \mathcal F_\beta(R)
 =\Tr(R\log R)
 +\beta\left[S_{\rm rad}(R)+2e_2(R)\right],
 \label{eq:v6-canonical-thermal-exterior-free-energy}
\end{equation}
где $S_{\rm rad}$ уже выведен после минимизации общей амплитуды моста.

На одноосной семье
\begin{equation}
 R(a)=\operatorname{diag}\left(a,\frac{1-a}{2},\frac{1-a}{2}\right)
\end{equation}
найден переход первого рода при
\begin{equation}
 \boxed{\beta_c=1.5426695409\ldots.}
 \label{eq:v6-canonical-cooling-critical-beta}
\end{equation}
В точке сосуществования упорядоченная фаза имеет спектр
\begin{equation}
 \Spec R_*
 =(0.9121666\ldots,0.0439167\ldots,0.0439167\ldots)
 \label{eq:v6-canonical-cooling-coexistence-spectrum}
\end{equation}
и общую сингулярную амплитуду $T_*=1.1011276\ldots$.

При $\beta<\beta_c$ глобален изотропный минимум $I_3/3$. При
$\beta>\beta_c$ глобальна полноранговая одноосная орбита
$SO(3)/O(2)=\mathbb{RP}^2$. Например, при $\beta=2$ выделенное
собственное значение равно $0.9760191\ldots$.

\begin{theorem}[Каноническая кристаллизация при охлаждении]
Дополнительный множитель два внешнему квадрату не требуется, если
отношение энергии к энтропии является термодинамической переменной.
Канонический родительский вес $2e_2(R)$ сам создаёт переход
$I_3/3\to\mathbb{RP}^2$ при конечной критической обратной температуре
$\beta_c$.
\end{theorem}

Это первый механизм Тома VI, буквально реализующий раннюю идею мира как
кристаллизации: высокотемпературная среда изотропна, охлаждение создаёт
ориентационные домены, а несовпавшие домены допускают проекторные дефекты.
Нематический изотропно-ориентационный переход действительно является
переходом первого рода в признанных моделях Майера--Заупе; см.
\path{doi:10.1103/PhysRevE.72.041703} и
\path{doi:10.1103/PhysRevE.72.041702}.

\section{Связь с существующим модулярным параметром}

Проект уже содержит семейство верных состояний
\begin{equation}
 \rho_\beta=\frac{e^{-\beta h_F}}{\Tr e^{-\beta h_F}},
 \qquad \beta>0,
\end{equation}
согласующееся с гипотезой теплового времени Конна--Ровелли; см.
\path{arXiv:gr-qc/9406019}. Поэтому появление обратной температуры не
является новой разновидностью поля или коэффициента связи.

Но существующий $h_F$ скалярен внутри семейных триплетов. Его модулярный
поток фиксирует проекторную ось и сохраняет собственное состояние. Кроме
того, проект не вычисляет абсолютное значение $\beta$ и не выводит путь
\begin{equation}
 \beta<\beta_c\longrightarrow\beta>\beta_c.
\end{equation}
Собственный модулярный поток не может охладить само породившее его
состояние. Для рождения доменов требуется семейство неравновесных состояний,
относительный модулярный коцикл, расширение или разрежение носителя либо
реляционные часы, задающие изменение эффективной температуры.

\begin{nogocriterion}[Запрет самопроизвольного охлаждения одним состоянием]
Наличие параметра $\beta$ в $\rho_\beta$ не доказывает, что система
пересекает $\beta_c$. Динамика охлаждения должна быть выведена отдельно и
не может совпадать с обратимым модулярным потоком фиксированного состояния.
\end{nogocriterion}

\section{Вердикт}

\begin{protocol}
\begin{enumerate}
 \item вклад тензорного квадрата одной ориентации:
 \textbf{$2e_2$};
 \item вклад смешанных ориентаций: \textbf{не зависящая от ранга константа};
 \item дополнительный множитель из Real-полуследа: \textbf{нет};
 \item ненормированный коэффициент четыре как нормированный проектор обмена:
 \textbf{нет};
 \item двойной $SO(3)$-казимир: \textbf{точно равен четырём};
 \item $SO(3)_L\times SO(3)_R$-лапласиан в текущем родителе:
 \textbf{не выведен};
 \item канонический внешний вес: \textbf{$2e_2(R)$};
 \item критическая обратная температура:
 \textbf{$\beta_c=1.5426695409\ldots$};
 \item высокотемпературная фаза: \textbf{изотропная};
 \item низкотемпературная фаза: \textbf{одноосная $\mathbb{RP}^2$};
 \item рабочий механизм кристаллизации: \textbf{получен статически};
 \item динамика пересечения $\beta_c$: \textbf{не выведена};
 \item следующая проверка: \textbf{модулярное охлаждение и проекторный переход};
 \item рождение материи: \textbf{существенно переоткрыто, но не закрыто}.
\end{enumerate}
\end{protocol}

Следующий тест должен связать термодинамический параметр
\eqref{eq:v6-canonical-thermal-exterior-free-energy} с уже построенным
модулярным состоянием и проверить, может ли относительный коцикл или
четырёхтактная подсистема часов породить монотонное пересечение
$\beta_c$ без внешнего времени.

Машинный сертификат сохранён в
\path{s2t_v6_tensor_square_relative_carrier_normalization_gate_results.json}.